\par\addvspace{9pt}\Needspace{4\baselineskip}\noindent{\bfseries\fontsize{12}{17}\selectfont 补充课程：18.433 组合优化}\par\nobreak\vspace{4pt}
官方课程名为 \textit{Combinatorial Optimization}，Santosh Vempala 主讲，2003 年秋季，本科课程；每周两次课，每次 1.5 小时。先修课为 18.06 或 18.700（线性代数）。课程以线性及凸规划、流与匹配、随机化和近似算法处理离散优化。本册仅吸收其中的图论、匹配、流、割及图优化题，不改编全部线性规划内容。

\begin{longtable}{p{0.12\textwidth}p{0.74\textwidth}}
\toprule 讲次 & 官网进度与关键事件\\\midrule\endhead
1--3 & 匹配理论；匈牙利算法；Edmonds 算法。\\
4--6 & 多面体组合学；匹配多面体两讲。\\
7--11 & 流与对偶；最大流（第 8 讲作业 A1 截止）；最小割；最小费用流；强多项式算法。\\
12--16 & 线性规划对偶；单纯形算法（第 13 讲 A2 截止）；第 14 讲考试 I；单纯形续讲；互补松弛与原始—对偶算法。\\
17--20 & 椭球法思想与细节；凸问题和组合问题的分离预言机。\\
21--25 & NP 完全性（第 21 讲 A3 截止）；两讲近似算法；松弛与舍入；第 25 讲考试 II，A4 截止。\\
26--27 & 项目回顾。\\\bottomrule
\end{longtable}
考核为考试 I、II 各占 25\%，项目研究 20\%，项目展示 6\%，作业 24\%。项目可个人或二人合作研究问题或实验评价算法，提交短报告并展示；应在第 5 次课后的次日确定项目。公开档案有四份作业；课程安排中的两场考试未公开试卷或答案，不能把考试 II 擅称为期末统考。

官网把下列资料列作补充参考，未指定唯一教材，也未给逐讲阅读页码、出版年或版本：Lovasz，\textit{Matching Theory}；Ahuja、Magnanti、Orlin，\textit{Network Flows}；Schrijver，\textit{Theory of Linear and Integer Programming}；Chvatal，\textit{Linear Programming}；Bertsimas、Tsitsiklis，\textit{Linear Optimization}；Cook、Cunningham、Pulleyblank、Schrijver，\textit{Combinatorial Optimization}；Papadimitriou、Steiglitz，同名 \textit{Combinatorial Optimization}。这里忠实保留官网书目缩写，不替官网补全缺项。
\sourceof{\href{https://ocw.mit.edu/courses/18-433-combinatorial-optimization-fall-2003/pages/syllabus/}{18.433 大纲}、\href{https://ocw.mit.edu/courses/18-433-combinatorial-optimization-fall-2003/pages/calendar/}{进度}及\href{https://ocw.mit.edu/courses/18-433-combinatorial-optimization-fall-2003/pages/readings/}{阅读书目}。官网按讲次列进度，未列具体上课日期。}

\par\addvspace{9pt}\Needspace{4\baselineskip}\noindent{\bfseries\fontsize{12}{17}\selectfont 补充课程：18.212 代数组合学}\par\nobreak\vspace{4pt}
官方课程名为 \textit{Algebraic Combinatorics}，Alexander Postnikov 主讲，Andrew Lin 整理讲义，2019 年春季，本科课程；每周三次课，每次一小时。先修课为 18.701（代数 I）或 18.703（现代代数）。大纲包括 Catalan 数与路径、对称群及排列统计、偏序与格、Young 表与 RSK、整数分拆、不交路径、生成树、停车函数、矩阵树、电网络、随机游走、染色多项式、群作用计数及多面体。本册选取图、树、相关路径和偏序题，不声称覆盖全部代数组合课程。
\begin{longtable}{p{0.12\textwidth}p{0.74\textwidth}}
\toprule 讲次 & 官网进度与关键事件\\\midrule\endhead
1--2 & Catalan 数、反射原理、路径与二叉树。\\
3--6 & 模式避免、Young 表、Schensted 对应、钩长和树型偏序。\\
7--14 & $q$ 计数、排列统计、Stirling 数、Eulerian 数；第 12 讲 PSet1 截止，第 13--14 讲讨论。\\
15--21 & 偏序、格、Sperner 与 Dilworth 定理、微分偏序、整数分拆。\\
22--25 & Cayley 树计数；第 23 讲 PSet2 截止，第 24--25 讲讨论。\\
26--32 & 矩阵树、超立方体树数、有向树、电网络、BEST 定理。\\
33--34 & 停车函数、树逆序统计、Bernoulli 数与 Catalan 数。\\
35--39 & 不交路径、平面分拆与铺砌、Schur 多项式；第 37 讲 PSet3 截止，第 38 讲讨论。\\\bottomrule
\end{longtable}
成绩以三份作业为依据，官网未给数值权重。PSet1 鼓励尽可能多做，但说明满分不要求全部提交，约六题可足够；这一说明不扩写为三份作业的统一规则。公开资料含五份参考解答文件及相应讨论讲义；部分题没有官方答案，个别官方解答只有部分论证。本册逐题说明所采用的参考答案及编者补充。

阅读书目为 Richard P. Stanley：\textit{Algebraic Combinatorics: Walks, Trees, Tableaux, and More}（Springer，2018；官网另列 2013 年在线稿），以及 \textit{Enumerative Combinatorics} 第 1 卷（课程列 1997）、第 2 卷（课程列 2001）。图论相关指定阅读包括第 6 讲 Knuth《计算机程序设计艺术》第 3 卷（1998）\S5.1.4 的根树钩长公式，第 17 讲前述代数组合书第 4 章的 Sperner 性质，第 18 讲计数组合第 1 卷 \S3.21 的微分偏序，第 23 讲代数组合书第 9 章附录的三种组合证明及 Egecioglu--Remmel 的 1986 年论文。仅列书目与指引，不复制受限书内题面。
\sourceof{\href{https://ocw.mit.edu/courses/18-212-algebraic-combinatorics-spring-2019/pages/syllabus/}{18.212 大纲}、\href{https://ocw.mit.edu/courses/18-212-algebraic-combinatorics-spring-2019/pages/calendar/}{进度}及\href{https://ocw.mit.edu/courses/18-212-algebraic-combinatorics-spring-2019/pages/readings/}{阅读}。日期未列出时只保留讲次；讨论 PDF 首页日期另见来源清单。}

\par\addvspace{9pt}\Needspace{4\baselineskip}\noindent{\bfseries\fontsize{12}{17}\selectfont 补充课程：18.217 图论与加法组合}\par\nobreak\vspace{4pt}
官方课程名为 \textit{Graph Theory and Additive Combinatorics}，Yufei Zhao 主讲，2019 年秋季，研究生课程；每周两次课，每次 1.5 小时。官网没有指定先修课程号，要求数学成熟度达到数学研究生一年级。课程研究禁子图、正则引理、拟随机图、图极限、Roth 定理、集合加法结构及和积问题。本册吸收禁子图和谱图论的部分主线及习题；正则性、图极限的其余题不冒称为非图论，加法组合部分不在本册范围。
\begin{longtable}{p{0.12\textwidth}p{0.74\textwidth}}
\toprule 讲次 & 官网进度与关键事件\\\midrule\endhead
1 & 图论与加法组合之间的桥梁。\\
2--5 & Mantel 与 Turán 定理；完全二部禁子图；代数构造；依赖随机选择。\\
6--10 & 图正则引理的陈述与证明；三角形消去；其他应用；诱导消去；超图消去与谱证明。第 6、10 讲分别作业 1、2 截止。\\
11--13 & 拟随机性；第二特征值；稀疏正则性与 Green--Tao 定理。第 13 讲作业 3 截止。\\
14--17 & 图极限导论；正则性与计数；紧致性与应用；子图密度不等式。第 17 讲作业 4 截止。\\
18--20 & Roth 定理：有限域上的 Fourier 证明、整数中的 Fourier 证明、多项式法与算术正则性。第 20 讲作业 5 截止。\\
21--25 & 集合加法结构：Freiman 定理、有限指数群与模型引理、Bogolyubov 引理与数的几何、Freiman 定理证明、加法能量与 Balog--Szemerédi--Gowers 定理。\\
26 & 和积问题与关联几何；作业 6 截止。\\\bottomrule
\end{longtable}
课程无考试。成绩取作业表现和写作表现两类中的较低者，边界成绩可参考课堂参与；写作作业包括课堂笔记和维基百科贡献。只做非星号题可达到 A$-$，A 或 A$+$ 需要解决较多星号题；不能把每一道星号题都称为开放问题。官网未公布数值权重。六次提交作业使用同一份题库，其“不提交”小问在本册中仍单独标明。

大纲和进度没有指定必读教材或逐周阅读书目。公开的完整讲义及八个专题章由学生依据课堂记录，Zhao 教授参与编辑，官网注明经许可使用。讲义中的参考文献不改标为课程指定阅读；2023 年作者图书只提供链接，不随本册再分发。
\sourceof{\href{https://ocw.mit.edu/courses/18-217-graph-theory-and-additive-combinatorics-fall-2019/pages/syllabus/}{18.217 大纲}、\href{https://ocw.mit.edu/courses/18-217-graph-theory-and-additive-combinatorics-fall-2019/pages/calendar/}{进度}及\href{https://ocw.mit.edu/courses/18-217-graph-theory-and-additive-combinatorics-fall-2019/pages/lecture-notes/}{讲义说明}。官网列 26 讲，没有具体上课日期。}

\par\addvspace{9pt}\Needspace{4\baselineskip}\noindent{\bfseries\fontsize{12}{17}\selectfont 补充课程：6.042J／18.062J 计算机科学数学}\par\nobreak\vspace{4pt}
官方课程名为 \textit{Mathematics for Computer Science}，Albert R. Meyer、Adam Chlipala 主讲，2015 年春季，本科课程；两课程号交叉开设。每周三次课，每次 1.5 小时，采用翻转课堂，课前阅读与在线反馈，课堂以六至八人团队解决问题。先修课为 18.01（一元微积分），涉及数列、级数、极限、微分与积分。原课程分数学基础、离散结构、离散概率三部分；本册仅使用图论及必要的图上随机游走题。
\begin{longtable}{p{0.12\textwidth}p{0.46\textwidth}p{0.28\textwidth}}
\toprule 课次 & 官网主题 & 教材阅读\\\midrule\endhead
1--5 & 证明导论、证明方法、良序原理、命题逻辑、量词与谓词。 & 1.1--1.9；2.1--2.3；3.1--3.6\\
6--8 & 集合、二元关系、归纳法。 & 4.1--4.5；5.1--5.3\\
9--11 & 状态机与不变量、递归定义、无限集合。 & 5.4；第6、7章\\
12--15 & 最大公因数、同余、数论中的 Euler 定理、RSA。 & 8.1--8.12\\
16 & 有向图：游走与路径。 & 9.1--9.4\\
17 & 有向无环图。 & 9.5\\
18 & 偏序与等价关系。 & 9.5--9.11\\
19 & 度数与同构。 & 11.1--11.4\\
20 & 染色与连通性。 & 11.7--11.9\\
21 & 树。 & 11.9--11.10\\
22 & 稳定匹配。 & 11.6\\
23--27 & 和与积、渐近、双射计数、重复与二项式、鸽巢与容斥。 & 13.1--13.5、13.7；14.1--14.2、14.4--14.8\\
28--34 & 离散概率、条件概率、独立性、随机变量、期望、偏差界、抽样与置信度。 & 16.1--16.5；17.1--17.5、17.7--17.8；18.1--18.5；19.1--19.5\\
35 & 随机游走与 PageRank。 & 20.2\\\bottomrule
\end{longtable}
教材为 Eric Lehman、F. Tom Leighton、Albert R. Meyer 合著的 \textit{Mathematics for Computer Science}（课程 2015 年版本）；教材作者与本期授课教师分别登记。上表合并的课次保持原有先后次序，具体每次课阅读范围见本册来源索引。

考核权重为课堂参与 25\%、在线反馈 5\%、作业 15\%、三场期中考试合计 30\%、期末考试 25\%。每场期中 80 分钟，可带一张双面资料纸；期末 180 分钟，可带两张双面资料纸。最低一次作业与最低三次课堂参与不计总评；课堂和在线题主要按参与计分，每份作业设计为至多三小时完成。公开档案有 12 份作业、35 份课堂题、三份期中和一份期末；这些 PDF 的官方解答不向一般公众开放，另有公开在线反馈题及解释，不能把两者混为一谈。2019 年 Open Learning Library 重发布的对应 PDF 与 2015 年归档逐个字节相同，重发布年不替代题目的编写年。
\sourceof{\href{https://ocw.mit.edu/courses/6-042j-mathematics-for-computer-science-spring-2015/pages/syllabus/}{6.042J 大纲}、\href{https://ocw.mit.edu/courses/6-042j-mathematics-for-computer-science-spring-2015/pages/readings/}{课次与阅读}及\href{https://openlearninglibrary.mit.edu/courses/course-v1:OCW+6.042J+2T2019/about}{Open Learning Library 重发布说明}。课程页面只给课次，不编造上课日历。}
